\begin{proof}[Proof of \cref{obj:lem:original-record-mixtures}]
\begin{enumerate}
\item Fix the joint certificate \((\varepsilon,M)\) from
\cref{obj:lem:calibrated-support}, and fix admissible sample sizes and
amplitudes. Since \(\exp(1),\exp(2)>0\) and \(M\ge1\),
\[
\kappa>0,\qquad C_{\mathrm{mix}}>0.
\]
Write
\[
\mathcal I_n=\{1,\ldots,n\},\qquad
\mathscr L=\{0,1\}^2,\qquad
\tau=\frac14\operatorname{count}_{\mathscr L}.
\]
Under the identification of an ordered record sample with its covariate
and label vectors, the reference measure is
\[
R=\lambda^{\otimes n}\otimes\tau^{\otimes n}.
\]
For a law \(P\) in the calibrated families, define its conditional label
density by
\[
\mathsf f_P(x,l)
=4P\{(A,Y)=l\mid X=x\},
\qquad x\in[0,1],\quad l\in\mathscr L.
\]
The support certificate supplies
\[
\frac12\le \mathsf f_P(x,l)\le2,
\qquad
\frac14\sum_{l\in\mathscr L}\mathsf f_P(x,l)=1
\]
throughout the prescribed signed amplitude domains.
% lean: CausalSmith.Stat.LogoddsLowsmoothFrontier.calib_constants_spec

\item We first bound the covariate-dependent component weight. For
\(\mathbf x=(x_i)_{i\in\mathcal I_n}\in[0,1]^n\), define
\[
j_k(x)=\min\{k-1,\lfloor kx\rfloor\},
\qquad
\mathscr G_{k,\mathbf x}:
\quad
i\sim i'
\ \Longleftrightarrow\
i\ne i'\ \text{ and }\ |j_k(x_i)-j_k(x_{i'})|\le1.
\]
Thus the cells are half-open, with \(x=1\) assigned to the last cell.
Let
\[
\mathscr V_{\mathbf x}
=\{\text{vertex sets of connected components of }\mathscr G_{k,\mathbf x}\},
\qquad
\mathscr B_{\mathbf x}=\mathscr V_{\mathbf x}\cup\{\varnothing\}.
\]
For every \(\mathcal V\in\mathscr B_{\mathbf x}\), put
\[
m_{\mathcal V}=|\mathcal V|,
\qquad
\mathsf w_{\mathcal V}
=
\begin{cases}
m_{\mathcal V}^4\,8^{m_{\mathcal V}},&m_{\mathcal V}\ge2,\\
0,&m_{\mathcal V}\le1,
\end{cases}
\]
and define
\[
\mathsf W_{\mathrm{occ}}(\mathbf x)
=\sum_{\mathcal V\in\mathscr B_{\mathbf x}}\mathsf w_{\mathcal V},
\qquad
\overline{\mathsf W}_{\mathrm{occ}}
=\int\mathsf W_{\mathrm{occ}}(\mathbf x)\,
d\lambda^{\otimes n}(\mathbf x).
\]
Counting components by their vertex sets gives the identity
\[
\mathsf W_{\mathrm{occ}}(\mathbf x)
=\sum_{m_{\mathrm{occ}}=2}^{n+1}
m_{\mathrm{occ}}^4\,8^{m_{\mathrm{occ}}}
\sum_{\substack{\mathcal I_{\mathrm{rec}}\subseteq\mathcal I_n\\
|\mathcal I_{\mathrm{rec}}|=m_{\mathrm{occ}}}}
\mathbf1_{\{\mathcal I_{\mathrm{rec}}\in\mathscr V_{\mathbf x}\}}.
\]
The term of size \(n+1\) is zero.

Here is the component probability estimate used to integrate this
identity. Fix a nonempty set \(\mathcal I_{\mathrm{rec}}\subseteq\mathcal I_n\)
and a root \(i_{\mathrm{root}}\in\mathcal I_{\mathrm{rec}}\). Define the set
of parent arrays by
\[
\mathscr P_{\mathrm{rec}}
=\left\{
\pi_{\mathrm{par}}:
\mathcal I_{\mathrm{rec}}\setminus\{i_{\mathrm{root}}\}
\longrightarrow\mathcal I_{\mathrm{rec}}
\right\}.
\]
An array is valid when there exists a height function satisfying
\[
\mathsf h_{\mathrm{par}}:\mathcal I_{\mathrm{rec}}\to\mathbb N_0,
\qquad
\mathsf h_{\mathrm{par}}(i_{\mathrm{root}})=0,
\qquad
\mathsf h_{\mathrm{par}}(\pi_{\mathrm{par}}(i))
<\mathsf h_{\mathrm{par}}(i)
\quad(i\ne i_{\mathrm{root}}).
\]
For a valid array, define its edge event by
\[
\mathscr A_{\pi_{\mathrm{par}}}
=
\left\{\mathbf x:
|j_k(x_i)-j_k(x_{\pi_{\mathrm{par}}(i)})|\le1
\text{ for every }i\ne i_{\mathrm{root}}\right\}.
\]
If \(\mathcal I_{\mathrm{rec}}\) is a component, choose each nonroot
vertex's parent one step closer to the root in the induced graph.
Graph distance then supplies the required height function. Consequently,
\[
\mathbf1_{\{\mathcal I_{\mathrm{rec}}\in\mathscr V_{\mathbf x}\}}
\le
\sum_{\substack{\pi_{\mathrm{par}}\in\mathscr P_{\mathrm{rec}}\\
\pi_{\mathrm{par}}\ \mathrm{valid}}}
\mathbf1_{\mathscr A_{\pi_{\mathrm{par}}}}(\mathbf x),
\qquad
|\mathscr P_{\mathrm{rec}}|
=|\mathcal I_{\mathrm{rec}}|^{|\mathcal I_{\mathrm{rec}}|-1}.
\]

For any fixed parent covariate \(x_{\mathrm{par}}\in[0,1]\),
\[
\left\{x_{\mathrm{child}}:
|j_k(x_{\mathrm{child}})-j_k(x_{\mathrm{par}})|\le1\right\}
\subseteq
\left[
\frac{j_k(x_{\mathrm{par}})-1}{k},
\frac{j_k(x_{\mathrm{par}})+2}{k}
\right]\cap[0,1].
\]
Indeed, a point in cell \(j_k(x_{\mathrm{child}})\) lies between that
cell's two endpoints, and the two cell numbers differ by at most one.
The displayed interval has length at most \(3/k\); hence
\[
\lambda\left\{x_{\mathrm{child}}:
|j_k(x_{\mathrm{child}})-j_k(x_{\mathrm{par}})|\le1\right\}
\le\frac3k.
\]
In a valid array, a nonroot vertex of maximal height has no children.
Integrating its covariate first removes one edge and costs at most
\(3/k\). Repeating this operation leaves the unrestricted root, whose
mass is one. Thus
\[
\lambda^{\otimes n}(\mathscr A_{\pi_{\mathrm{par}}})
\le
\left(\frac3k\right)^{|\mathcal I_{\mathrm{rec}}|-1},
\]
including the one-vertex case, where the edge product is empty. Summing
over parent arrays yields
\[
\int
\mathbf1_{\{\mathcal I_{\mathrm{rec}}\in\mathscr V_{\mathbf x}\}}\,
d\lambda^{\otimes n}(\mathbf x)
\le
|\mathcal I_{\mathrm{rec}}|^{|\mathcal I_{\mathrm{rec}}|-1}
\left(\frac3k\right)^{|\mathcal I_{\mathrm{rec}}|-1}.
\]
Therefore
\[
\overline{\mathsf W}_{\mathrm{occ}}
\le
\sum_{m_{\mathrm{occ}}=2}^{n+1}
m_{\mathrm{occ}}^4\,8^{m_{\mathrm{occ}}}
\binom n{m_{\mathrm{occ}}}
m_{\mathrm{occ}}^{m_{\mathrm{occ}}-1}
\left(\frac3k\right)^{m_{\mathrm{occ}}-1}.
\]

To evaluate this finite envelope, use the nonnegative exponential series:
for every integer \(m_{\mathrm{occ}}\ge1\),
\[
\frac{m_{\mathrm{occ}}^{m_{\mathrm{occ}}}}{m_{\mathrm{occ}}!}
\le \exp(m_{\mathrm{occ}})
=\exp(1)^{m_{\mathrm{occ}}},
\qquad
\frac{m_{\mathrm{occ}}^{m_{\mathrm{occ}}-1}}
{m_{\mathrm{occ}}!}
\le\frac{\exp(1)^{m_{\mathrm{occ}}}}{m_{\mathrm{occ}}}.
\]
Also,
\[
\binom n{m_{\mathrm{occ}}}
\le\frac{n^{m_{\mathrm{occ}}}}{m_{\mathrm{occ}}!},
\]
because each factor in the falling factorial is at most \(n\), with
the binomial coefficient equal to zero when \(m_{\mathrm{occ}}>n\).
Set
\[
q_{\mathrm{occ}}=24\exp(1)\frac nk.
\]
The occupancy hypothesis gives \(0\le q_{\mathrm{occ}}\le1/4\), and the
preceding two estimates imply, term by term,
\[
\begin{aligned}
&m_{\mathrm{occ}}^4\,8^{m_{\mathrm{occ}}}
\binom n{m_{\mathrm{occ}}}
m_{\mathrm{occ}}^{m_{\mathrm{occ}}-1}
\left(\frac3k\right)^{m_{\mathrm{occ}}-1}\\
&\hspace{2em}\le
8\exp(1)n\,m_{\mathrm{occ}}^3
q_{\mathrm{occ}}^{m_{\mathrm{occ}}-1}.
\end{aligned}
\]

For completeness, the finite cubic sum has the required numerical bound
by an elementary induction. Define
\[
\mathsf A_{\mathrm{tail}}(h)
=2h^3+12h^2+30h+32,
\qquad h\in\mathbb N_0.
\]
For \(h,N_{\mathrm{tail}}\in\mathbb N_0\),
\[
\sum_{h_{\mathrm{sum}}=0}^{N_{\mathrm{tail}}-1}
(h+h_{\mathrm{sum}}+2)^3q_{\mathrm{occ}}^{h_{\mathrm{sum}}}
\le \mathsf A_{\mathrm{tail}}(h).
\]
The empty sum satisfies this inequality. For the induction step, split
off the first term and apply the induction hypothesis at \(h+1\);
the resulting bound follows from
\[
\begin{aligned}
(h+2)^3+q_{\mathrm{occ}}\mathsf A_{\mathrm{tail}}(h+1)
&\le (h+2)^3+\frac14\mathsf A_{\mathrm{tail}}(h+1),\\
\mathsf A_{\mathrm{tail}}(h)
-\left((h+2)^3+\frac14\mathsf A_{\mathrm{tail}}(h+1)\right)
&=\frac12h^3+\frac32h^2+3h+5\ge0.
\end{aligned}
\]
At \(h=0\), multiplication by \(q_{\mathrm{occ}}\) gives
\[
\sum_{h_{\mathrm{sum}}=0}^{n-1}
(h_{\mathrm{sum}}+2)^3q_{\mathrm{occ}}^{h_{\mathrm{sum}}+1}
\le32q_{\mathrm{occ}}.
\]
Combining the estimates,
\[
\overline{\mathsf W}_{\mathrm{occ}}
\le8\exp(1)n\,32q_{\mathrm{occ}}
=96\exp(2)\,8^2\frac{n^2}{k}.
\]
In particular, for every \(\omega_{\mathrm{amp}}\in\mathbb R\),
\[
M^4\omega_{\mathrm{amp}}^2\overline{\mathsf W}_{\mathrm{occ}}
\le
C_{\mathrm{mix}}\frac{n^2}{k}\omega_{\mathrm{amp}}^2.
\]
The weights are integrable finite sums of measurable component
indicators.
% lean: CausalSmith.Stat.LogoddsLowsmoothFrontier.mixture_weighted_component_integral_bound

\item We now prove the mixed inequality, first treating its degenerate
cases. If \(n=1\), signed singleton matching in
\cref{obj:lem:calibrated-support} gives
\[
\mathbb E_\sigma\mathsf f_{P_{0,\sigma}}(x,l)
=\mathbb E_\sigma\mathsf f_{P_{1,\sigma}}(x,l),
\qquad
\mathcal H^2(Q_0,Q_1)=0.
\]
If \(\eta=0\), the two four-cell tables in
\cref{obj:def:continuous-calibrated-priors} agree for each sign vector:
the propensity is \(1/2\), the effect is zero, and the covariance
adjustment at zero effect is zero. Thus \(Q_0=Q_1\).
If \(\zeta=0\), reversing all signs gives instead the draw-wise identity
\[
P_{0,\sigma}(\eta,0)=P_{1,-\sigma}(\eta,0).
\]
Here the outcome margin is \(1/2\), the covariance adjustment is zero,
and sign reversal exchanges the two propensity margins. Since
\(\sigma\mapsto-\sigma\) permutes the uniform finite prior, again
\(Q_0=Q_1\). In both amplitude cases the Hellinger distance and its
claimed upper bound are zero. Henceforth, for the mixed comparison,
take \(n\ge2\), \(\eta\ne0\), and \(\zeta\ne0\).
% lean: CausalSmith.Stat.LogoddsLowsmoothFrontier.mixed_originalRecordHellinger_one
% lean: CausalSmith.Stat.LogoddsLowsmoothFrontier.mixedMixture_zero_left
% lean: CausalSmith.Stat.LogoddsLowsmoothFrontier.mixedMixture_zero_right

\item Fix the full covariate vector \(\mathbf x\). We describe the
conditional component decomposition explicitly. For each nonempty
\(\mathcal V\in\mathscr B_{\mathbf x}\), and for the empty block, define
\[
\mathscr S_{\mathcal V}
=\bigcup_{i\in\mathcal V}\{j_k(x_i),j_k(x_i)+1\},
\qquad
\mathscr S_{\varnothing}
=\{0,\ldots,k\}\setminus
\bigcup_{\mathcal V\in\mathscr V_{\mathbf x}}\mathscr S_{\mathcal V}.
\]
Distinct nonempty components have disjoint endpoint-sign sets: sharing
an endpoint would make the corresponding cell numbers equal or
adjacent. These sets therefore partition the \(k+1\) endpoint signs.
For an assignment
\(\sigma_{\mathcal V}\in\{-1,1\}^{\mathscr S_{\mathcal V}}\), use the
extension
\[
\widehat\sigma_{\mathcal V}(s)
=
\begin{cases}
\sigma_{\mathcal V}(s),&s\in\mathscr S_{\mathcal V},\\
-1,&s\notin\mathscr S_{\mathcal V}.
\end{cases}
\]
Every record in \(\mathcal V\) depends only on the assigned signs.
Each assignment has exactly \(2^{k+1-|\mathscr S_{\mathcal V}|}\)
completions to a full sign vector, so restricting the uniform prior
preserves each block average.

For \(b\in\{0,1\}\), define the full and component conditional densities
by
\[
\begin{aligned}
\mathsf L_b^{\mathrm{mix}}(\mathbf x,\boldsymbol l;\eta,\zeta)
&=\mathbb E_\sigma
\prod_{i\in\mathcal I_n}
\mathsf f_{P_{b,\sigma}}(x_i,l_i),\\
\mathsf L_{b,\mathcal V}^{\mathrm{mix}}
(\boldsymbol l_{\mathcal V};\eta,\zeta)
&=\mathbb E_{\sigma_{\mathcal V}}
\prod_{i\in\mathcal V}
\mathsf f_{P_{b,\widehat\sigma_{\mathcal V}}}(x_i,l_i),
\qquad
\boldsymbol l_{\mathcal V}\in\mathscr L^{\mathcal V}.
\end{aligned}
\]
The covariates remain fixed in the component notation. The empty
component density is one. Independence of the disjoint sign groups
and finite distributivity give
\[
\mathsf L_b^{\mathrm{mix}}(\mathbf x,\boldsymbol l;\eta,\zeta)
=
\prod_{\mathcal V\in\mathscr B_{\mathbf x}}
\mathsf L_{b,\mathcal V}^{\mathrm{mix}}
(\boldsymbol l|_{\mathcal V};\eta,\zeta).
\]

To make the Hellinger tensor bound transparent, define component label
masses and affinities by
\[
\mathsf p_{b,\mathcal V}^{\mathrm{mix}}(\boldsymbol l_{\mathcal V})
=4^{-m_{\mathcal V}}
\mathsf L_{b,\mathcal V}^{\mathrm{mix}}
(\boldsymbol l_{\mathcal V};\eta,\zeta),
\qquad
\mathfrak a_{\mathcal V}^{\mathrm{mix}}
=
\sum_{\boldsymbol l_{\mathcal V}\in\mathscr L^{\mathcal V}}
\sqrt{\mathsf p_{0,\mathcal V}^{\mathrm{mix}}
(\boldsymbol l_{\mathcal V})
\mathsf p_{1,\mathcal V}^{\mathrm{mix}}
(\boldsymbol l_{\mathcal V})}.
\]
Normalization of each conditional four-cell table yields
\[
\begin{aligned}
\sum_{\boldsymbol l_{\mathcal V}}
\mathsf p_{b,\mathcal V}^{\mathrm{mix}}(\boldsymbol l_{\mathcal V})
&=
\mathbb E_{\sigma_{\mathcal V}}
\prod_{i\in\mathcal V}
\left(\frac14\sum_{l_i\in\mathscr L}
\mathsf f_{P_{b,\widehat\sigma_{\mathcal V}}}(x_i,l_i)\right)
=1,\\
\sum_{\boldsymbol l_{\mathcal V}}
\left(\sqrt{\mathsf p_{0,\mathcal V}^{\mathrm{mix}}}
-\sqrt{\mathsf p_{1,\mathcal V}^{\mathrm{mix}}}\right)^2
&=2(1-\mathfrak a_{\mathcal V}^{\mathrm{mix}}).
\end{aligned}
\]
All masses are nonnegative. The last equality, together with
nonnegativity of the squared sum, shows
\(0\le\mathfrak a_{\mathcal V}^{\mathrm{mix}}\le1\).
Regrouping labels by their component restrictions factors the full
affinity into the product of these affinities. The elementary
inequality
\[
1-\prod_{\mathcal V}\mathfrak a_{\mathcal V}^{\mathrm{mix}}
\le
\sum_{\mathcal V}(1-\mathfrak a_{\mathcal V}^{\mathrm{mix}})
\]
follows by induction from
\(1-aa'=(1-a)+a(1-a')\le(1-a)+(1-a')\) for
\(a,a'\in[0,1]\).

Define the conditional discrepancy by
\[
\mathsf h_{\mathrm{mix}}(\mathbf x)
=
\int_{\mathscr L^{\mathcal I_n}}
\left(
\sqrt{\mathsf L_0^{\mathrm{mix}}(\mathbf x,\boldsymbol l;\eta,\zeta)}
-\sqrt{\mathsf L_1^{\mathrm{mix}}(\mathbf x,\boldsymbol l;\eta,\zeta)}
\right)^2\,d\tau^{\otimes n}(\boldsymbol l).
\]
Since each label vector has reference mass \(4^{-n}\), scaling the
densities by \(4^{-n}\) turns this integral exactly into the squared
root discrepancy of the corresponding label masses. The affinity
calculation therefore proves
\[
\mathsf h_{\mathrm{mix}}(\mathbf x)
\le
\sum_{\mathcal V\in\mathscr B_{\mathbf x}}
4^{-m_{\mathcal V}}
\sum_{\boldsymbol l_{\mathcal V}\in\mathscr L^{\mathcal V}}
\left(
\sqrt{\mathsf L_{0,\mathcal V}^{\mathrm{mix}}}
-\sqrt{\mathsf L_{1,\mathcal V}^{\mathrm{mix}}}
\right)^2.
\]
Finally, the common uniform covariate marginal and the displayed
product densities give, by integration over covariates and labels,
\[
\mathcal H^2(Q_0,Q_1)
=\int\mathsf h_{\mathrm{mix}}(\mathbf x)\,
d\lambda^{\otimes n}(\mathbf x).
\]
% lean: CausalSmith.Stat.LogoddsLowsmoothFrontier.mixed_conditional_label_hellinger_le_components
% lean: CausalSmith.Stat.LogoddsLowsmoothFrontier.mixed_originalRecordHellinger_conditional

\item We bound each mixed component's pointwise square-root discrepancy.
If \(m_{\mathcal V}=0\), both densities are one. If
\(m_{\mathcal V}=1\), restriction of the uniform prior and signed
singleton matching in \cref{obj:lem:calibrated-support} give equality
of the two densities. Thus, for \(m_{\mathcal V}\le1\),
\[
\left(
\sqrt{\mathsf L_{0,\mathcal V}^{\mathrm{mix}}}
-\sqrt{\mathsf L_{1,\mathcal V}^{\mathrm{mix}}}
\right)^2=0=M^4(\eta\zeta)^2\mathsf w_{\mathcal V}.
\]

Suppose \(m_{\mathcal V}\ge2\), and fix its labels. Define, throughout
the signed amplitude square,
\[
\mathsf D_{\mathcal V}^{\mathrm{mix}}(\eta_0,\zeta_0)
=
\mathsf L_{0,\mathcal V}^{\mathrm{mix}}
(\boldsymbol l_{\mathcal V};\eta_0,\zeta_0)
-
\mathsf L_{1,\mathcal V}^{\mathrm{mix}}
(\boldsymbol l_{\mathcal V};\eta_0,\zeta_0),
\qquad
|\eta_0|,|\zeta_0|\le\varepsilon.
\]
For a fixed sign assignment, put
\[
\mathsf f_{b,i}^{\mathrm{mix}}(\eta_0,\zeta_0)
=\mathsf f_{P_{b,\widehat\sigma_{\mathcal V}}}(x_i,l_i),
\qquad
\mathsf F_{b,\mathcal J}^{\mathrm{mix}}(\eta_0,\zeta_0)
=\prod_{i\in\mathcal J}\mathsf f_{b,i}^{\mathrm{mix}}(\eta_0,\zeta_0),
\quad \mathcal J\subseteq\mathcal V.
\]
Here each law on the right is evaluated at \((\eta_0,\zeta_0)\).
The derivative and common-neighborhood assertions of
\cref{obj:lem:calibrated-support} imply
\[
|\mathsf f_{b,i}^{\mathrm{mix}}|\le2,\qquad
|\partial_{\eta_0}\mathsf f_{b,i}^{\mathrm{mix}}|,
|\partial_{\zeta_0}\mathsf f_{b,i}^{\mathrm{mix}}|,
|\partial_{\zeta_0}\partial_{\eta_0}\mathsf f_{b,i}^{\mathrm{mix}}|
\le M.
\]
The partial derivative bounds follow by evaluating the total
derivatives on unit coordinate vectors.

Product differentiation gives, with \(m_{\mathcal J}=|\mathcal J|\),
\[
\begin{aligned}
|\mathsf F_{b,\mathcal J}^{\mathrm{mix}}|&\le2^{m_{\mathcal J}},\\
|\partial_{\eta_0}\mathsf F_{b,\mathcal J}^{\mathrm{mix}}|,
|\partial_{\zeta_0}\mathsf F_{b,\mathcal J}^{\mathrm{mix}}|
&\le M m_{\mathcal J}2^{m_{\mathcal J}},\\
|\partial_{\zeta_0}\partial_{\eta_0}
\mathsf F_{b,\mathcal J}^{\mathrm{mix}}|
&\le M^2m_{\mathcal J}^2 2^{m_{\mathcal J}}.
\end{aligned}
\]
These estimates can be proved together by induction on
\(\mathcal J\). For the empty product its value is one and all its
derivatives vanish. On adjoining \(i_{\mathrm{new}}\), the first
derivative bound is controlled by
\[
M2^{m_{\mathcal J}}+
2M m_{\mathcal J}2^{m_{\mathcal J}}
\le
M(m_{\mathcal J}+1)2^{m_{\mathcal J}+1}.
\]
The mixed product rule has four terms:
\[
\begin{aligned}
\partial_{\zeta_0}\partial_{\eta_0}
\left(\mathsf f_{b,i_{\mathrm{new}}}^{\mathrm{mix}}
\mathsf F_{b,\mathcal J}^{\mathrm{mix}}\right)
={}&
(\partial_{\zeta_0}\partial_{\eta_0}
\mathsf f_{b,i_{\mathrm{new}}}^{\mathrm{mix}})
\mathsf F_{b,\mathcal J}^{\mathrm{mix}}
+
(\partial_{\eta_0}\mathsf f_{b,i_{\mathrm{new}}}^{\mathrm{mix}})
(\partial_{\zeta_0}\mathsf F_{b,\mathcal J}^{\mathrm{mix}})\\
&+
(\partial_{\zeta_0}\mathsf f_{b,i_{\mathrm{new}}}^{\mathrm{mix}})
(\partial_{\eta_0}\mathsf F_{b,\mathcal J}^{\mathrm{mix}})
+
\mathsf f_{b,i_{\mathrm{new}}}^{\mathrm{mix}}
(\partial_{\zeta_0}\partial_{\eta_0}
\mathsf F_{b,\mathcal J}^{\mathrm{mix}}).
\end{aligned}
\]
Their absolute sum is at most
\[
\left(M+2M^2m_{\mathcal J}
+2M^2m_{\mathcal J}^2\right)2^{m_{\mathcal J}}
\le
M^2(m_{\mathcal J}+1)^2 2^{m_{\mathcal J}+1},
\]
using \(M\ge1\). Differentiation commutes with the finite sign
averages, so subtraction of the two averages yields
\[
|\partial_{\zeta_0}\partial_{\eta_0}
\mathsf D_{\mathcal V}^{\mathrm{mix}}(\eta_0,\zeta_0)|
\le
2M^2m_{\mathcal V}^2 2^{m_{\mathcal V}}.
\]
The one-record density floor also gives
\[
\mathsf L_{0,\mathcal V}^{\mathrm{mix}},
\mathsf L_{1,\mathcal V}^{\mathrm{mix}}
\ge2^{-m_{\mathcal V}}.
\]

The same table identities used in the amplitude cases, now applied
to a block average, imply
\[
\mathsf D_{\mathcal V}^{\mathrm{mix}}(0,\zeta_0)=0,\qquad
\mathsf D_{\mathcal V}^{\mathrm{mix}}(\eta_0,0)=0,\qquad
\partial_{\eta_0}\mathsf D_{\mathcal V}^{\mathrm{mix}}(\eta_0,0)=0.
\]
For the second identity, reverse the component signs, a permutation
of their uniform prior. The two closed segments joining zero to
\(\eta\) and \(\zeta\) lie inside the certificate's signed square.
The mean value bound along the \(\zeta_0\)-segment therefore gives
\[
|\partial_{\eta_0}\mathsf D_{\mathcal V}^{\mathrm{mix}}(\eta_0,\zeta)|
\le
2M^2m_{\mathcal V}^2 2^{m_{\mathcal V}}|\zeta|
\]
for every \(\eta_0\) between zero and \(\eta\). Applying the mean value
bound along that \(\eta_0\)-segment next gives
\[
|\mathsf D_{\mathcal V}^{\mathrm{mix}}(\eta,\zeta)|
\le
2M^2m_{\mathcal V}^2 2^{m_{\mathcal V}}|\eta\zeta|.
\]
Both estimates apply to either orientation of the segments.

The positive density floors justify the square-root identity and
bound
\[
\begin{aligned}
\left(
\sqrt{\mathsf L_{0,\mathcal V}^{\mathrm{mix}}}
-\sqrt{\mathsf L_{1,\mathcal V}^{\mathrm{mix}}}
\right)^2
&=
\frac{(\mathsf D_{\mathcal V}^{\mathrm{mix}}(\eta,\zeta))^2}
{\left(
\sqrt{\mathsf L_{0,\mathcal V}^{\mathrm{mix}}}
+\sqrt{\mathsf L_{1,\mathcal V}^{\mathrm{mix}}}
\right)^2}\\
&\le
\frac{4M^4m_{\mathcal V}^4 4^{m_{\mathcal V}}(\eta\zeta)^2}
{4\,2^{-m_{\mathcal V}}}\\
&=M^4(\eta\zeta)^2\mathsf w_{\mathcal V}.
\end{aligned}
\]
This proves the pointwise component bound for every block, including
the unused-sign block.
% lean: CausalSmith.Stat.LogoddsLowsmoothFrontier.mixedDensity_small_block_matching
% lean: CausalSmith.Stat.LogoddsLowsmoothFrontier.mixed_component_hellinger_from_support

\item Averaging the pointwise bound over a component's labels preserves
it, since there are \(4^{m_{\mathcal V}}\) label vectors:
\[
4^{-m_{\mathcal V}}
\sum_{\boldsymbol l_{\mathcal V}}
\left(
\sqrt{\mathsf L_{0,\mathcal V}^{\mathrm{mix}}}
-\sqrt{\mathsf L_{1,\mathcal V}^{\mathrm{mix}}}
\right)^2
\le M^4(\eta\zeta)^2\mathsf w_{\mathcal V}.
\]
Summing over components and integrating over the original covariates
now yields
\[
\begin{aligned}
\mathcal H^2(Q_0,Q_1)
&\le
\int M^4(\eta\zeta)^2\mathsf W_{\mathrm{occ}}(\mathbf x)\,
d\lambda^{\otimes n}(\mathbf x)\\
&=M^4(\eta\zeta)^2\overline{\mathsf W}_{\mathrm{occ}}\\
&\le C_{\mathrm{mix}}\frac{n^2}{k}\eta^2\zeta^2.
\end{aligned}
\]
The conditional discrepancy is nonnegative, and the integrable
component weight supplies the integrable upper bound.
% lean: CausalSmith.Stat.LogoddsLowsmoothFrontier.conditionalComponent_mass_hellinger_le_density_bound
% lean: CausalSmith.Stat.LogoddsLowsmoothFrontier.original_record_mixtures

\item Fix \(t\in[0,1/4]\) for the fair comparison. If \(n=1\), signed
singleton matching in \cref{obj:lem:calibrated-support} gives
\[
\mathbb E_\sigma\mathsf f_{P_{t,\sigma}}(x,l)
=\mathsf f_{P_{*,t,\delta}}(x,l),
\qquad
\mathcal H^2(Q_t,P_{*,t,\delta})=0.
\]
If \(\delta=0\), the formulas in
\cref{obj:def:calibration-maps,obj:def:continuous-calibrated-priors}
give \(\mathsf T(t,0)=t\); the random and comparator tables then agree
for every sign draw. Thus
\[
Q_t=P_{*,t,0}^{\otimes n},
\qquad
\mathcal H^2(Q_t,P_{*,t,0}^{\otimes n})=0.
\]
These cases satisfy the desired inequality. Take henceforth \(n\ge2\)
and \(\delta\ne0\).
% lean: CausalSmith.Stat.LogoddsLowsmoothFrontier.fair_originalRecordHellinger_one
% lean: CausalSmith.Stat.LogoddsLowsmoothFrontier.fairMixture_zero_amplitude

\item First establish a reduction for any auxiliary amplitude
\(\delta_{\mathrm{aux}}\in[0,\varepsilon]\). With the covariates fixed,
define
\[
\begin{aligned}
\mathsf L_t^{\mathrm{fair}}(\mathbf x,\boldsymbol l;\delta_{\mathrm{aux}})
&=\mathbb E_\sigma\prod_{i\in\mathcal I_n}
\mathsf f_{P_{t,\sigma}}(x_i,l_i),\\
\mathsf L_*^{\mathrm{fair}}(\mathbf x,\boldsymbol l;\delta_{\mathrm{aux}})
&=\prod_{i\in\mathcal I_n}
\mathsf f_{P_{*,t,\delta_{\mathrm{aux}}}}(x_i,l_i),\\
\mathsf L_{t,\mathcal V}^{\mathrm{fair}}
(\boldsymbol l_{\mathcal V};\delta_{\mathrm{aux}})
&=\mathbb E_{\sigma_{\mathcal V}}\prod_{i\in\mathcal V}
\mathsf f_{P_{t,\widehat\sigma_{\mathcal V}}}(x_i,l_i),\\
\mathsf L_{*,\mathcal V}^{\mathrm{fair}}
(\boldsymbol l_{\mathcal V};\delta_{\mathrm{aux}})
&=\prod_{i\in\mathcal V}
\mathsf f_{P_{*,t,\delta_{\mathrm{aux}}}}(x_i,l_i).
\end{aligned}
\]
The random laws in these displays are evaluated at
\(\delta_{\mathrm{aux}}\). Averaging a sign-independent product is
exactly that product:
\[
\mathbb E_\sigma
\prod_{i\in\mathcal I_n}
\mathsf f_{P_{*,t,\delta_{\mathrm{aux}}}}(x_i,l_i)
=
\prod_{i\in\mathcal I_n}
\mathsf f_{P_{*,t,\delta_{\mathrm{aux}}}}(x_i,l_i).
\]
Restriction of the random prior to each component and finite
factorization consequently give
\[
\mathsf L_t^{\mathrm{fair}}
=\prod_{\mathcal V\in\mathscr B_{\mathbf x}}
\mathsf L_{t,\mathcal V}^{\mathrm{fair}},
\qquad
\mathsf L_*^{\mathrm{fair}}
=\prod_{\mathcal V\in\mathscr B_{\mathbf x}}
\mathsf L_{*,\mathcal V}^{\mathrm{fair}}.
\]
For each of the two component densities, summing its scaled label
mass gives
\[
4^{-m_{\mathcal V}}
\sum_{\boldsymbol l_{\mathcal V}}
\mathsf L_{t,\mathcal V}^{\mathrm{fair}}
=
4^{-m_{\mathcal V}}
\sum_{\boldsymbol l_{\mathcal V}}
\mathsf L_{*,\mathcal V}^{\mathrm{fair}}
=1.
\]
For the random density this follows by commuting the finite sign
average with the label sum; in both cases the label sum factors into
normalized one-record sums.

Define the component affinities and conditional discrepancy by
\[
\begin{aligned}
\mathfrak a_{\mathcal V}^{\mathrm{fair}}
&=
4^{-m_{\mathcal V}}\sum_{\boldsymbol l_{\mathcal V}}
\sqrt{\mathsf L_{t,\mathcal V}^{\mathrm{fair}}
\mathsf L_{*,\mathcal V}^{\mathrm{fair}}},\\
\mathsf h_{\mathrm{fair}}(\mathbf x;\delta_{\mathrm{aux}})
&=
\int
\left(
\sqrt{\mathsf L_t^{\mathrm{fair}}}
-\sqrt{\mathsf L_*^{\mathrm{fair}}}
\right)^2\,d\tau^{\otimes n}.
\end{aligned}
\]
Expanding the squared difference and using normalization gives
\(0\le\mathfrak a_{\mathcal V}^{\mathrm{fair}}\le1\) and
\[
\begin{aligned}
\mathsf h_{\mathrm{fair}}(\mathbf x;\delta_{\mathrm{aux}})
&=2\left(1-\prod_{\mathcal V}
\mathfrak a_{\mathcal V}^{\mathrm{fair}}\right)\\
&\le
2\sum_{\mathcal V}(1-\mathfrak a_{\mathcal V}^{\mathrm{fair}})\\
&=
\sum_{\mathcal V\in\mathscr B_{\mathbf x}}
4^{-m_{\mathcal V}}\sum_{\boldsymbol l_{\mathcal V}}
\left(
\sqrt{\mathsf L_{t,\mathcal V}^{\mathrm{fair}}}
-\sqrt{\mathsf L_{*,\mathcal V}^{\mathrm{fair}}}
\right)^2.
\end{aligned}
\]
Here the affinity factors by finite distributivity, and the inequality
is the finite-product inequality for numbers in \([0,1]\).
The common covariate marginal also gives
\[
\mathcal H^2(Q_t,P_{*,t,\delta_{\mathrm{aux}}}^{\otimes n})
=
\int\mathsf h_{\mathrm{fair}}(\mathbf x;\delta_{\mathrm{aux}})
\,d\lambda^{\otimes n}(\mathbf x),
\]
where \(Q_t\) in this display is evaluated at the auxiliary amplitude.
% lean: CausalSmith.Stat.LogoddsLowsmoothFrontier.fair_conditional_label_hellinger_le_components
% lean: CausalSmith.Stat.LogoddsLowsmoothFrontier.fair_originalRecordHellinger_conditional

\item For \(m_{\mathcal V}\le1\), the two fair component densities agree:
the empty products are one, and the singleton identity follows from
signed singleton matching in \cref{obj:lem:calibrated-support} and
restriction of the uniform sign prior. Hence their square-root
discrepancy is zero.

For \(m_{\mathcal V}\ge2\), fix the component labels and define, on the
whole signed interval,
\[
\mathsf D_{\mathcal V}^{\mathrm{fair}}(\delta_0)
=
\mathsf L_{t,\mathcal V}^{\mathrm{fair}}
(\boldsymbol l_{\mathcal V};\delta_0)
-
\mathsf L_{*,\mathcal V}^{\mathrm{fair}}
(\boldsymbol l_{\mathcal V};\delta_0),
\qquad |\delta_0|\le\varepsilon.
\]
For the derivative calculation, define the one-record factors and
their partial products by
\[
\begin{aligned}
\mathsf f_{t,i}^{\mathrm{fair}}(\delta_0)
&=\mathsf f_{P_{t,\widehat\sigma_{\mathcal V}}}(x_i,l_i),&
\mathsf f_{*,i}^{\mathrm{fair}}(\delta_0)
&=\mathsf f_{P_{*,t,\delta_0}}(x_i,l_i),\\
\mathsf F_{t,\mathcal J}^{\mathrm{fair}}(\delta_0)
&=\prod_{i\in\mathcal J}\mathsf f_{t,i}^{\mathrm{fair}}(\delta_0),&
\mathsf F_{*,\mathcal J}^{\mathrm{fair}}(\delta_0)
&=\prod_{i\in\mathcal J}\mathsf f_{*,i}^{\mathrm{fair}}(\delta_0),
\qquad \mathcal J\subseteq\mathcal V.
\end{aligned}
\]
The random law on the right is evaluated at \(\delta_0\). By
\cref{obj:lem:calibrated-support}, both kinds of factors are twice
continuously differentiable and satisfy
\[
|\mathsf f_{t,i}^{\mathrm{fair}}|,
|\mathsf f_{*,i}^{\mathrm{fair}}|\le2,\qquad
|(\mathsf f_{t,i}^{\mathrm{fair}})'|,
|(\mathsf f_{*,i}^{\mathrm{fair}})'|,
|(\mathsf f_{t,i}^{\mathrm{fair}})''|,
|(\mathsf f_{*,i}^{\mathrm{fair}})''|\le M.
\]
For either kind of partial product, the product rule and induction
give
\[
|\mathsf F|\le2^{m_{\mathcal J}},\qquad
|\mathsf F'|\le M m_{\mathcal J}2^{m_{\mathcal J}},\qquad
|\mathsf F''|\le M^2m_{\mathcal J}^2 2^{m_{\mathcal J}}.
\]
To see the second-derivative induction explicitly, adjoining a
factor gives
\[
(\mathsf f\,\mathsf F)''
=\mathsf f\,\mathsf F''+2\mathsf f'\mathsf F'
+\mathsf f''\mathsf F,
\]
whose absolute value is at most
\[
\left(2M^2m_{\mathcal J}^2+2M^2m_{\mathcal J}+M\right)
2^{m_{\mathcal J}}
\le
M^2(m_{\mathcal J}+1)^2 2^{m_{\mathcal J}+1}.
\]
Here \(\mathsf f\) denotes the adjoined one-record factor and
\(\mathsf F\) the product over \(\mathcal J\), of either kind. The
first-derivative induction uses
\[
M2^{m_{\mathcal J}}+
2M m_{\mathcal J}2^{m_{\mathcal J}}
\le M(m_{\mathcal J}+1)2^{m_{\mathcal J}+1}.
\]
The empty product again has value one and zero derivatives.
Finite averaging and subtraction therefore imply, for
\(0\le\delta_0\le\delta_{\mathrm{aux}}\),
\[
|(\mathsf D_{\mathcal V}^{\mathrm{fair}})''(\delta_0)|
\le2M^2m_{\mathcal V}^2 2^{m_{\mathcal V}}.
\]
The density floors give
\[
\mathsf L_{t,\mathcal V}^{\mathrm{fair}},
\mathsf L_{*,\mathcal V}^{\mathrm{fair}}
\ge2^{-m_{\mathcal V}}.
\]

At zero amplitude the random and comparator tables coincide, so
\[
\mathsf D_{\mathcal V}^{\mathrm{fair}}(0)=0.
\]
The fair root is even in its amplitude by
\cref{obj:lem:exact-calibrations}; the comparator effect in
\cref{obj:def:calibration-maps} depends on the amplitude through its
square. Reversing the component signs consequently proves
\[
\mathsf D_{\mathcal V}^{\mathrm{fair}}(-\delta_0)
=\mathsf D_{\mathcal V}^{\mathrm{fair}}(\delta_0).
\]
Differentiating this equality at zero, using the supplied smoothness,
gives
\[
(\mathsf D_{\mathcal V}^{\mathrm{fair}})'(0)
=-(\mathsf D_{\mathcal V}^{\mathrm{fair}})'(0)=0.
\]
If \(\delta_{\mathrm{aux}}=0\), the difference is already zero. If
\(\delta_{\mathrm{aux}}>0\), Taylor's theorem with Lagrange remainder
provides \(\delta_{\mathrm{mid}}\in(0,\delta_{\mathrm{aux}})\) such that
\[
\mathsf D_{\mathcal V}^{\mathrm{fair}}(\delta_{\mathrm{aux}})
=
\frac{\delta_{\mathrm{aux}}^2}{2}
(\mathsf D_{\mathcal V}^{\mathrm{fair}})''(\delta_{\mathrm{mid}}).
\]
Thus in either case,
\[
|\mathsf D_{\mathcal V}^{\mathrm{fair}}(\delta_{\mathrm{aux}})|
\le
M^2m_{\mathcal V}^2 2^{m_{\mathcal V}}\delta_{\mathrm{aux}}^2
\le
2M^2m_{\mathcal V}^2 2^{m_{\mathcal V}}\delta_{\mathrm{aux}}^2.
\]
Using this last bound and the two positive floors,
\[
\begin{aligned}
\left(
\sqrt{\mathsf L_{t,\mathcal V}^{\mathrm{fair}}}
-\sqrt{\mathsf L_{*,\mathcal V}^{\mathrm{fair}}}
\right)^2
&=
\frac{(\mathsf D_{\mathcal V}^{\mathrm{fair}}
(\delta_{\mathrm{aux}}))^2}
{\left(
\sqrt{\mathsf L_{t,\mathcal V}^{\mathrm{fair}}}
+\sqrt{\mathsf L_{*,\mathcal V}^{\mathrm{fair}}}
\right)^2}\\
&\le
\frac{4M^4m_{\mathcal V}^4 4^{m_{\mathcal V}}
\delta_{\mathrm{aux}}^4}{4\,2^{-m_{\mathcal V}}}\\
&=M^4\delta_{\mathrm{aux}}^4\mathsf w_{\mathcal V}.
\end{aligned}
\]
Together with the small-component cases, this estimate holds for every
\(\mathcal V\in\mathscr B_{\mathbf x}\).
% lean: CausalSmith.Stat.LogoddsLowsmoothFrontier.fairDensity_small_block_matching
% lean: CausalSmith.Stat.LogoddsLowsmoothFrontier.fair_component_hellinger_from_support

\item Summing the fair pointwise bound over component labels with their
reference masses, then summing over components, gives
\[
\mathsf h_{\mathrm{fair}}(\mathbf x;\delta_{\mathrm{aux}})
\le
M^4(\delta_{\mathrm{aux}}^2)^2
\mathsf W_{\mathrm{occ}}(\mathbf x).
\]
Integration over the original covariates therefore establishes the
nonnegative-amplitude reduction
\[
\mathcal H^2(Q_t,P_{*,t,\delta_{\mathrm{aux}}}^{\otimes n})
\le
M^4(\delta_{\mathrm{aux}}^2)^2
\overline{\mathsf W}_{\mathrm{occ}},
\qquad 0\le\delta_{\mathrm{aux}}\le\varepsilon.
\]
As before, the conditional discrepancy is nonnegative and the
component weight is integrable.
% lean: CausalSmith.Stat.LogoddsLowsmoothFrontier.original_record_mixtures

\item Finally, restore the signed amplitude. The fair-root evenness and
the square dependence of the comparator effect give, throughout the
signed certificate interval,
\[
P_{t,\sigma}(-\delta)=P_{t,-\sigma}(\delta),
\qquad
P_{*,t,-\delta}=P_{*,t,\delta}.
\]
The notation in the first identity displays the amplitude argument of
the random fair law. Reindexing the finite sign average by
\(\sigma\mapsto-\sigma\) shows that the fair sample mixture is even
as well. Thus, if \(\delta\ge0\) use \(|\delta|=\delta\), and if
\(\delta<0\) use \(|\delta|=-\delta\), to obtain
\[
\mathcal H^2(Q_t,P_{*,t,\delta}^{\otimes n})
=
\left.
\mathcal H^2(Q_t,P_{*,t,\delta_{\mathrm{aux}}}^{\otimes n})
\right|_{\delta_{\mathrm{aux}}=|\delta|}.
\]
The mixture on the right is evaluated at \(|\delta|\).
Since \(0\le|\delta|\le\varepsilon\), the nonnegative-amplitude
reduction applies, and \(|\delta|^2=\delta^2\) gives
\[
\begin{aligned}
\mathcal H^2(Q_t,P_{*,t,\delta}^{\otimes n})
&\le M^4(\delta^2)^2\overline{\mathsf W}_{\mathrm{occ}}\\
&\le C_{\mathrm{mix}}\frac{n^2}{k}(\delta^2)^2\\
&=C_{\mathrm{mix}}\frac{n^2}{k}\delta^4.
\end{aligned}
\]
Every estimate used the same fixed certificate and was uniform over
\(t\in[0,1/4]\), proving both asserted bounds.
% lean: CausalSmith.Stat.LogoddsLowsmoothFrontier.fairMixture_even
% lean: CausalSmith.Stat.LogoddsLowsmoothFrontier.fairComparator_even
% lean: CausalSmith.Stat.LogoddsLowsmoothFrontier.original_record_mixtures
\end{enumerate}
\end{proof}
